\documentclass[11pt]{article}

\usepackage[margin=0.85in]{geometry}
\usepackage{amsmath,amssymb}
\usepackage{enumitem}
\usepackage{xcolor}
\usepackage{tcolorbox}
\usepackage{fancyhdr}
\usepackage{lastpage}
\usepackage{hyperref}
\usepackage{microtype}

\definecolor{uwpurple}{HTML}{4B2E83}
\definecolor{uwgold}{HTML}{B7A57A}
\definecolor{softpurple}{HTML}{F4F1FA}
\definecolor{softgold}{HTML}{FBF7EA}
\definecolor{deepgold}{HTML}{8A6F2A}

\tcbuselibrary{breakable,skins}

\newtcolorbox{solutionbox}[1][]{
    breakable,
    enhanced,
    colback=softgold,
    colframe=deepgold,
    coltitle=white,
    colbacktitle=uwpurple,
    fonttitle=\bfseries,
    title={Solution #1},
    boxrule=0.85pt,
    arc=2mm,
    left=2.5mm,
    right=2.5mm,
    top=1.5mm,
    bottom=1.5mm,
    attach boxed title to top left={xshift=2.5mm,yshift=-2mm},
    boxed title style={
        colback=uwpurple,
        colframe=uwpurple,
        arc=1.5mm,
        boxrule=0pt,
        left=1.6mm,
        right=1.6mm,
        top=0.6mm,
        bottom=0.6mm
    }
}

\newtcolorbox{finalanswerbox}[1][]{
    colback=white,
    colframe=uwpurple,
    boxrule=0.75pt,
    arc=2mm,
    left=2mm,
    right=2mm,
    top=1mm,
    bottom=1mm,
    title={\textbf{Final Answer #1}},
    coltitle=uwpurple,
    colbacktitle=softpurple,
    fonttitle=\bfseries
}

\hypersetup{
    colorlinks=true,
    linkcolor=uwpurple,
    urlcolor=uwpurple
}

\pagestyle{fancy}
\fancyhf{}
\lhead{\textbf{MATH 394: Probability I}}
\rhead{\textbf{Homework 4 Solutions}}
\cfoot{\thepage\ of \pageref{LastPage}}

\setlength{\parindent}{0pt}
\setlength{\parskip}{0.45em}

\newcommand{\course}{MATH 394: Probability I}
\newcommand{\instructor}{Arman Jahangiri}
\newcommand{\department}{Department of Mathematics}
\newcommand{\university}{University of Washington}

\newtcolorbox{problemheader}{
    colback=softpurple,
    colframe=uwpurple,
    boxrule=0.8pt,
    arc=2mm,
    left=2mm,
    right=2mm,
    top=1mm,
    bottom=1mm
}

\newcommand{\source}[1]{%
    \vspace{0.15cm}
    {\small\textit{Source: Blitzstein and Hwang, \textit{Introduction to Probability}, Chapter 3, Problem #1.}}
}

\begin{document}

\begin{titlepage}
\centering
\vspace*{1.2cm}

{\Large \textbf{\university}}\\
{\large \department}

\vspace{1.5cm}

{\Huge \color{uwpurple}\textbf{Homework 4 Solutions}}\\[0.25cm]
{\Large \textbf{\course}}\\[0.25cm]
{\large Instructor: Arman Jahangiri}

\vspace{1cm}

\begin{tcolorbox}[
    width=0.78\textwidth,
    colback=softpurple,
    colframe=uwpurple,
    boxrule=0.9pt,
    arc=3mm
]
\centering
\textbf{Submission:} Single PDF on Gradescope

\textbf{Deadline:} 10:00 PM Pacific Time on the listed due date
\end{tcolorbox}

\vspace{0.8cm}

\begin{tcolorbox}[
    width=0.86\textwidth,
    colback=white,
    colframe=uwgold,
    boxrule=0.8pt,
    arc=3mm
]
This homework is worth \textbf{50 points}.

Across the quarter, there are \textbf{eight homework assignments}, worth \textbf{400 points total}. Homework assignments together account for \textbf{40\%} of the final course grade. Further course policies can be seen in the \href{https://armanjg.github.io/MATH-394-Probability-1-v2/}{\textbf{MATH 394 syllabus}.}
\end{tcolorbox}

\vfill
\end{titlepage}

\section*{Homework Policy}

\subsection*{Submission and Deadline}

Please:

\begin{itemize}[leftmargin=*]
    \item Submit homework as a single PDF on Gradescope by \textbf{10:00 PM (Pacific Time)} on the listed due date.
    \item Begin each problem on a new page, and ensure that pages are correctly tagged in Gradescope.
    \item Write solutions clearly and legibly.
    \item Typesetting solutions in LaTeX is encouraged, but not required.
\end{itemize}

\subsection*{Late Days}

Each student is allotted \textbf{six late days} for the quarter. A late day extends a homework deadline by up to 24 hours without penalty. For example, submitting an assignment anytime between 10:01 PM on the due date and 10:00 PM the following day counts as one late day.

The following rules apply:
\begin{itemize}[leftmargin=*]
   \item No explanation is required to use late days.
   \item At most \textbf{two late days} may be used on any single homework assignment.
   \item Late days may not be used on the final homework assignment or on any assignment for which solutions have already been released.
   \item It is the student's responsibility to keep track of how many late days have been used during the quarter.
\end{itemize}

Once all late days have been exhausted, additional late submissions will incur a penalty of \textbf{10\% per day}, up to a maximum deduction of \textbf{50\%}. Assignments submitted more than five days late, or after solutions have been released, will not be accepted.

\subsection*{Submission Issues and Technical Difficulties}

If a serious technical issue prevents a timely Gradescope submission, students may temporarily submit their assignment by email to the instructor at \href{mailto:armanjg@uw.edu}{armanjg@uw.edu}. In such cases:

\begin{itemize}[leftmargin=*]
    \item The submission must consist of a \textbf{single PDF file}. The email subject line must follow the format:
    \[
    \texttt{HW\# Submission - FULL NAME - STUDENT ID}.
    \]
    \item Once the Gradescope issue has been resolved, the student must upload the \emph{exact same file} to Gradescope as soon as possible.
    \item The student should clearly indicate in the Gradescope submission that the assignment was previously submitted by email before the deadline.
\end{itemize}

Email submissions are intended only for genuine technical emergencies and should not be used as a substitute for timely Gradescope submission.

\newpage

% ------------------------------------------------------------

\begin{problemheader}
\textbf{Problem 1. }
\end{problemheader}
\source{1}

People are arriving at a party one at a time. While waiting for more people to arrive, they entertain themselves by comparing their birthdays.

Let $X$ be the number of people needed to obtain a birthday match; that is, before person $X$ arrives, no two people have the same birthday, but when person $X$ arrives there is a match.

Find the PMF of $X$.

\begin{solutionbox}
Assume the usual birthday model with 365 equally likely birthdays and no leap day. The smallest possible value of $X$ is 2, since one person alone cannot create a match. The largest possible value is 366, since after 365 people with distinct birthdays, the next person must match one of them.

For $X=k$, the first $k-1$ people must all have distinct birthdays, and the $k$th person must match one of those $k-1$ birthdays. The probability that the first $k-1$ birthdays are distinct is
\[
\frac{365\cdot 364\cdots (365-k+2)}{365^{k-1}}.
\]
Once this has happened, the $k$th person has exactly $k-1$ favorable birthdays out of 365. Therefore, for $k=2,3,\ldots,366$,
\[
P(X=k)
=
\frac{365\cdot 364\cdots (365-k+2)}{365^{k-1}}\cdot \frac{k-1}{365}.
\]

 

\begin{finalanswerbox}
\[
P(X=k)=\frac{365\cdot 364\cdots (365-k+2)}{365^{k-1}}\cdot \frac{k-1}{365},
\qquad k=2,3,\ldots,366.
\]
\[
P(X=k)=0 \qquad \text{otherwise.}
\]
\end{finalanswerbox}
\end{solutionbox}
\newpage

% ------------------------------------------------------------

\begin{problemheader}
\textbf{Problem 2. }
\end{problemheader}
\source{3}

Let $X$ be a random variable with CDF $F$, and let
\[
Y=\mu+\sigma X,
\]
where $\mu$ and $\sigma$ are real numbers with $\sigma>0$.

The random variable $Y$ is called a \emph{location--scale transformation} of $X$.

Find the CDF of $Y$ in terms of $F$.

\begin{solutionbox}
Let $F_Y$ denote the CDF of $Y$. For any real number $y$,
\[
F_Y(y)=P(Y\le y).
\]
Since $Y=\mu+\sigma X$ and $\sigma>0$, the inequality may be rearranged without reversing its direction:
\[
\mu+\sigma X\le y
\qquad \Longleftrightarrow \qquad
X\le \frac{y-\mu}{\sigma}.
\]
Thus
\[
F_Y(y)
=
P\left(X\le \frac{y-\mu}{\sigma}\right)
=
F\left(\frac{y-\mu}{\sigma}\right).
\]
 
\end{solutionbox}
\newpage

% ------------------------------------------------------------

\begin{problemheader}
\textbf{Problem 3.}
\end{problemheader}
\source{4}

Let $n$ be a positive integer and define
\[
F(x)=\frac{\lfloor x\rfloor}{n}
\]
for $0\le x\le n$, where $\lfloor x\rfloor$ is the greatest integer less than or equal to $x$.

Also define
\[
F(x)=0 \quad \text{for } x<0,
\]
and
\[
F(x)=1 \quad \text{for } x>n.
\]

\begin{enumerate}[label=(\alph*), leftmargin=*]
    \item Show that $F$ is a valid CDF.

    \item Find the PMF corresponding to $F$.
\end{enumerate}

\begin{solutionbox}
\textbf{(a)} The function $F$ is nondecreasing because $\lfloor x\rfloor$ is nondecreasing. It also has the correct limiting behavior. For $x<0$, $F(x)=0$, so
\[
\lim_{x\to -\infty}F(x)=0.
\]
For $x>n$, $F(x)=1$, so
\[
\lim_{x\to \infty}F(x)=1.
\]
The only possible jumps occur at the integers. At an integer $j$ with $1\le j\le n$, the left-hand limit is
\[
F(j^-)=\frac{j-1}{n},
\]
while
\[
F(j)=\frac{j}{n}.
\]
Thus the function is right-continuous at each jump. Away from the integers, it is constant on intervals and therefore right-continuous there as well. Hence $F$ has the required properties of a CDF.
 

\textbf{(b)} The PMF is obtained from the jumps of the CDF. At each integer $j=1,2,\ldots,n$,
\[
P(X=j)=F(j)-F(j^-)=\frac{j}{n}-\frac{j-1}{n}=\frac{1}{n}.
\]
There are no jumps elsewhere, so the probability is zero outside $\{1,2,\ldots,n\}$.
 
\end{solutionbox}
\newpage

% ------------------------------------------------------------

\begin{problemheader}
\textbf{Problem 4. }
\end{problemheader}
\source{5}

\begin{enumerate}[label=(\alph*), leftmargin=*]
    \item Show that
    \[
    p(n)=\left(\frac12\right)^{n+1},
    \qquad n=0,1,2,\dots,
    \]
    is a valid PMF for a discrete random variable.

    \item Find the CDF of a random variable with the PMF from part (a).
\end{enumerate}

\begin{solutionbox}
\textbf{(a)} Each value $p(n)$ is nonnegative. It remains to check that the probabilities add to 1. Using the geometric series,
\[
\sum_{n=0}^{\infty}\left(\frac12\right)^{n+1}
=
\frac12+\frac14+\frac18+\cdots
=
1.
\]
Therefore $p(n)$ is a valid PMF on the nonnegative integers.
 

\textbf{(b)} Let $X$ have this PMF. If $x<0$, then no possible value of $X$ is at most $x$, so $F(x)=0$. If $x\ge 0$, then
\[
F(x)=P(X\le x)
=
\sum_{n=0}^{\lfloor x\rfloor}\left(\frac12\right)^{n+1}.
\]
The finite geometric sum gives
\[
F(x)=1-\left(\frac12\right)^{\lfloor x\rfloor+1},
\qquad x\ge 0.
\]

\begin{finalanswerbox}[for Part (b)]
\[
F(x)=
\begin{cases}
0, & x<0,\\[1mm]
1-\left(\dfrac12\right)^{\lfloor x\rfloor+1}, & x\ge 0.
\end{cases}
\]
\end{finalanswerbox}
\end{solutionbox}
\newpage

% ------------------------------------------------------------

\begin{problemheader}
\textbf{Problem 5.}
\end{problemheader}
\source{8}

There are 100 prizes, with one worth \$1, one worth \$2, \dots, and one worth \$100. There are 100 boxes, each containing one of the prizes.

You get 5 prizes by picking random boxes one at a time, without replacement.

Find the PMF of how much your most valuable prize is worth, as a simple expression in terms of binomial coefficients.

\begin{solutionbox}
Let $M$ be the value of the most valuable prize among the 5 prizes selected. The sample space consists of all 5-element subsets of the 100 prizes, so there are
\[
\binom{100}{5}
\]
equally likely selections.

For the maximum value to be $m$, the prize worth $m$ dollars must be selected, and the other 4 selected prizes must all come from the prizes worth $1,2,\ldots,m-1$. This can be done in
\[
\binom{m-1}{4}
\]
ways. This is possible only when $m\ge 5$.

Therefore
\[
P(M=m)=\frac{\binom{m-1}{4}}{\binom{100}{5}},
\qquad m=5,6,\ldots,100.
\]
\[
P(M=m)=0 \qquad \text{otherwise.}
\]

\end{solutionbox}
\newpage

% ------------------------------------------------------------

\begin{problemheader}
\textbf{Problem 6.}
\end{problemheader}
\source{9}

Let $F_1$ and $F_2$ be CDFs, let $0<p<1$, and define
\[
F(x)=pF_1(x)+(1-p)F_2(x)
\]
for all $x$.

\begin{enumerate}[label=(\alph*), leftmargin=*]
    \item Show directly that $F$ has the properties of a valid CDF. The distribution defined by $F$ is called a \emph{mixture} of the distributions defined by $F_1$ and $F_2$.

    \item Consider creating a random variable in the following way. Flip a coin with probability $p$ of Heads. If the coin lands Heads, generate a random variable according to $F_1$; if the coin lands Tails, generate a random variable according to $F_2$. Show that the random variable obtained in this way has CDF $F$.
\end{enumerate}

\begin{solutionbox}
\textbf{(a)} Since $F_1$ and $F_2$ are nondecreasing and $p$ and $1-p$ are positive weights, their weighted average is also nondecreasing. Also,
\[
\lim_{x\to -\infty}F(x)
=
p\lim_{x\to -\infty}F_1(x)+(1-p)\lim_{x\to -\infty}F_2(x)
=
0,
\]
and
\[
\lim_{x\to \infty}F(x)
=
p\lim_{x\to \infty}F_1(x)+(1-p)\lim_{x\to \infty}F_2(x)
=
p+(1-p)=1.
\]
Finally, because $F_1$ and $F_2$ are right-continuous, for every $x$,
\[
\lim_{t\downarrow x}F(t)
=
p\lim_{t\downarrow x}F_1(t)+(1-p)\lim_{t\downarrow x}F_2(t)
=
pF_1(x)+(1-p)F_2(x)
=
F(x).
\]
Thus $F$ is a valid CDF.

\begin{finalanswerbox}[for Part (a)]
$F=pF_1+(1-p)F_2$ is nondecreasing, right-continuous, tends to $0$ at $-\infty$, and tends to $1$ at $\infty$. Hence it is a valid CDF.
\end{finalanswerbox}

\textbf{(b)} Let $C$ be the event that the coin lands Heads. The random variable is generated from $F_1$ on $C$ and from $F_2$ on $C^c$. Therefore, by the law of total probability,
\[
P(X\le x)
=
P(X\le x\mid C)P(C)+P(X\le x\mid C^c)P(C^c).
\]
Since $P(C)=p$, $P(C^c)=1-p$, $P(X\le x\mid C)=F_1(x)$, and $P(X\le x\mid C^c)=F_2(x)$, we get
\[
P(X\le x)=pF_1(x)+(1-p)F_2(x)=F(x).
\]
So the random variable constructed by this procedure has CDF $F$.

\begin{finalanswerbox}[for Part (b)]
\[
P(X\le x)=pF_1(x)+(1-p)F_2(x)=F(x).
\]
\end{finalanswerbox}
\end{solutionbox}
\newpage

% ------------------------------------------------------------

\begin{problemheader}
\textbf{Problem 7. }
\end{problemheader}
\source{22}

There are two coins, one with probability $p_1$ of Heads and the other with probability $p_2$ of Heads. One of the coins is randomly chosen, with equal probabilities for the two coins. It is then flipped $n\ge2$ times.

Let $X$ be the number of times it lands Heads.

\begin{enumerate}[label=(\alph*), leftmargin=*]
    \item Find the PMF of $X$.

    \item What is the distribution of $X$ if $p_1=p_2$?

    \item Give an intuitive explanation of why $X$ is not Binomial for $p_1\ne p_2$. You may assume that $n$ is large for your explanation.
\end{enumerate}

\begin{solutionbox}
\textbf{(a)} Conditional on choosing coin 1, the number of Heads has distribution $\operatorname{Bin}(n,p_1)$. Conditional on choosing coin 2, the number of Heads has distribution $\operatorname{Bin}(n,p_2)$. Since the two coins are chosen with equal probabilities,
\[
P(X=k)
=
\frac12\binom{n}{k}p_1^k(1-p_1)^{n-k}
+
\frac12\binom{n}{k}p_2^k(1-p_2)^{n-k},
\]
for $k=0,1,\ldots,n$.
 

\textbf{(b)} If $p_1=p_2=p$, then the two possible coins have the same Heads probability. The mixture collapses to the ordinary Binomial distribution,
\[
P(X=k)=\binom{n}{k}p^k(1-p)^{n-k}.
\]
Thus
\[
X\sim \operatorname{Bin}(n,p).
\]
 

\textbf{(c)} When $p_1\ne p_2$, the flips are not independent in the unconditional sense. Seeing many Heads early makes it more likely that the chosen coin is the one with the larger Heads probability, and that changes our prediction for later flips. For large $n$, the number of Heads tends to cluster around $np_1$ if coin 1 was chosen and around $np_2$ if coin 2 was chosen. This is a mixture of two Binomial patterns rather than one Binomial pattern with a single fixed success probability.
 
\end{solutionbox}
\newpage

% ------------------------------------------------------------

\begin{problemheader}
\textbf{Problem 8. }
\end{problemheader}
\source{30}

A certain company has $n+m$ employees, consisting of $n$ women and $m$ men. The company is deciding which employees to promote.

\begin{enumerate}[label=(\alph*), leftmargin=*]
    \item Suppose for this part that the company decides to promote $t$ employees, where
    \[
    1\le t\le n+m,
    \]
    by choosing $t$ random employees, with equal probabilities for each set of $t$ employees.

    What is the distribution of the number of women who get promoted?

    \item Now suppose that instead of having a predetermined number of promotions to give, the company decides independently for each employee, promoting the employee with probability $p$.

    Find the distributions of:
    \begin{itemize}
        \item the number of women who are promoted,
        \item the number of women who are not promoted,
        \item the number of employees who are promoted.
    \end{itemize}

    \item In the set-up from part (b), find the conditional distribution of the number of women who are promoted, given that exactly $t$ employees are promoted.
\end{enumerate}

\begin{solutionbox}
\textbf{(a)} Let $W$ be the number of women promoted. Since the company chooses $t$ employees from $n+m$ employees without replacement, $W$ has a Hypergeometric distribution. To have $W=k$, the company must choose $k$ of the $n$ women and $t-k$ of the $m$ men. Thus
\[
P(W=k)=\frac{\binom{n}{k}\binom{m}{t-k}}{\binom{n+m}{t}},
\]
where $k$ must satisfy
\[
\max(0,t-m)\le k\le \min(n,t).
\]
 

\textbf{(b)} Now each employee is promoted independently with probability $p$.

The number of women promoted is the number of successes among $n$ independent promotion decisions for women, so
\[
W_p\sim \operatorname{Bin}(n,p).
\]
The number of women not promoted is the number of failures among those same $n$ women, so
\[
W_{np}\sim \operatorname{Bin}(n,1-p).
\]
The total number of employees promoted is the number of successes among all $n+m$ employees, so
\[
T\sim \operatorname{Bin}(n+m,p).
\]

\begin{finalanswerbox}[for Part (b)]
\[
W_p\sim \operatorname{Bin}(n,p),
\qquad
W_{np}\sim \operatorname{Bin}(n,1-p),
\qquad
T\sim \operatorname{Bin}(n+m,p).
\]
\end{finalanswerbox}

\textbf{(c)} Given that exactly $t$ employees are promoted, all sets of $t$ promoted employees are equally likely. This brings us back to the same sampling-without-replacement structure as part (a). Therefore the conditional distribution of the number of women promoted is Hypergeometric:
\[
P(W=k\mid T=t)=\frac{\binom{n}{k}\binom{m}{t-k}}{\binom{n+m}{t}},
\]
for the same possible values of $k$, i.e.,  $\ \max(0,t-m)\le k\le \min(n,t).$

 
\end{solutionbox}

\end{document}
